\subsection{Construções em \textbf{Struct}}
\label{subsec:struct-constructions}

\subsubsection{Produto}

\begin{proposition}
O produto de $(d_1, \varphi_1, n_1)$ e $(d_2, \varphi_2, n_2)$ em \textbf{Struct} é o triplo
\[
    (d, \varphi, n_1 \cdot n_2),
\]
onde os elementos são indexados por pares $(i,j) \in \{1,\ldots,n_1\} \times \{1,\ldots,n_2\}$ (codificados como $(i-1)\cdot n_2 + j$), com
\[
    \varphi(i,j) = (\varphi_1(i), \varphi_2(j)), \qquad d(i,j) = (d_1(i), d_2(j)).
\]
\end{proposition}

\begin{lstlisting}[style=matlab, caption={Produto de dois objetos de Struct}]
function P = struct_product(S1, S2)
    P.n   = S1.n * S2.n;
    P.phi = zeros(1, P.n);
    P.d   = zeros(1, P.n, 'like', 1i);
    for i = 1:S1.n
        for j = 1:S2.n
            k        = (i-1)*S2.n + j;
            phi1_i   = S1.phi(i); phi2_j = S2.phi(j);
            P.phi(k) = (phi1_i-1)*S2.n + phi2_j;
            P.d(k)   = S1.d(i) + 1i*S2.d(j); % codificacao como complexo
        end
    end
end
\end{lstlisting}

\subsubsection{Coproduto}

O coproduto $(d_1, \varphi_1, n_1) \sqcup (d_2, \varphi_2, n_2)$ tem $n_1 + n_2$ elementos, com $\varphi$ e $d$ definidos por blocos:
\[
    \varphi(i) = \begin{cases} \varphi_1(i) & i \leq n_1 \\ n_1 + \varphi_2(i - n_1) & i > n_1 \end{cases}, \qquad
    d(i) = \begin{cases} d_1(i) & i \leq n_1 \\ d_2(i - n_1) & i > n_1 \end{cases}.
\]

\begin{lstlisting}[style=matlab, caption={Coproduto de dois objetos de Struct}]
function C = struct_coproduct(S1, S2)
    C.n   = S1.n + S2.n;
    C.phi = [S1.phi, S1.n + S2.phi];
    C.d   = [S1.d,   S2.d];
end
\end{lstlisting}

\subsubsection{Pullback}

Dado $S_1 \xrightarrow{f} S_3 \xleftarrow{g} S_2$, o pullback é:
\[
    \{(i,j) \in \{1,\ldots,n_1\} \times \{1,\ldots,n_2\} : f(i) = g(j)\},
\]
com endomapa $\varphi_{PB}(i,j) = (\varphi_1(i), \varphi_2(j))$ e $d_{PB}(i,j) = d_1(i)$ (que é igual a $d_2(j)$ pois $f$ e $g$ preservam $d$).

\begin{lstlisting}[style=matlab, caption={Pullback em Struct}]
function [PB, pi1, pi2] = struct_pullback(S1, S2, S3, f, g)
    pairs = [];
    for i = 1:S1.n
        for j = 1:S2.n
            if f(i) == g(j)
                pairs(end+1, :) = [i, j];
            end
        end
    end
    m    = size(pairs, 1);
    PB.n = m;
    PB.d = zeros(1, m, 'like', 1i);
    PB.phi = zeros(1, m);
    for k = 1:m
        i = pairs(k,1); j = pairs(k,2);
        PB.d(k) = S1.d(i);
        fi = S1.phi(i); fj = S2.phi(j);
        % Encontrar o indice do par (phi1(i), phi2(j)) em pairs
        idx = find(pairs(:,1)==fi & pairs(:,2)==fj, 1);
        PB.phi(k) = idx;
    end
    pi1 = pairs(:,1)';
    pi2 = pairs(:,2)';
end
\end{lstlisting}

\subsubsection{Pushout}

O pushout de $S_1 \xleftarrow{f} S_0 \xrightarrow{g} S_2$ é o coproduto $S_1 \sqcup S_2$ quocientado pela menor congruência compatível com $\varphi$ que identifica $f(k) \sim n_1 + g(k)$ para todo $k \in \{1,\ldots,n_0\}$.

\begin{lstlisting}[style=matlab, caption={Pushout em Struct (via union-find)}]
function [PO, q1, q2] = struct_pushout(S0, S1, S2, f, g)
    % Coproduto base
    C = struct_coproduct(S1, S2);
    % Union-find para as identificacoes f(k) ~ S1.n + g(k)
    parent = 1:C.n;
    function r = find_root(x)
        while parent(x) ~= x, x = parent(x); end
        r = x;
    end
    for k = 1:S0.n
        a = find_root(f(k));
        b = find_root(S1.n + g(k));
        parent(a) = b;
    end
    % Construir classes de equivalencia
    classes = zeros(1, C.n);
    rep     = 0;
    for x = 1:C.n
        r = find_root(x);
        if classes(r) == 0
            rep = rep + 1;
            classes(r) = rep;
        end
        classes(x) = classes(find_root(x));
    end
    PO.n   = rep;
    PO.phi = zeros(1, rep);
    PO.d   = zeros(1, rep, 'like', 1i);
    for x = 1:C.n
        cx = classes(x);
        PO.phi(cx) = classes(C.phi(x));
        PO.d(cx)   = C.d(x);
    end
    q1 = classes(1:S1.n);
    q2 = classes(S1.n+1:end);
end
\end{lstlisting}

\subsubsection{Equalizador}

O equalizador de $f, g: S_1 \rightrightarrows S_2$ é o subobjeto de $S_1$ induzido por $E = \{i : f(i) = g(i)\}$, desde que $E$ seja fechado para $\varphi_1$ (i.e., $\varphi_1(E) \subseteq E$).

\begin{lstlisting}[style=matlab, caption={Equalizador em Struct}]
function [E, inc] = struct_equalizer(S1, f, g)
    mask = (f == g);
    % Fechar para phi1
    changed = true;
    while changed
        changed = false;
        for i = 1:S1.n
            if mask(i) && ~mask(S1.phi(i))
                mask(S1.phi(i)) = false;
                mask(i) = false;
                changed = true;
            end
        end
    end
    inc  = find(mask);
    E.n  = numel(inc);
    E.d  = S1.d(inc);
    % Reindexar phi
    E.phi = zeros(1, E.n);
    for k = 1:E.n
        E.phi(k) = find(inc == S1.phi(inc(k)));
    end
end
\end{lstlisting}

\subsubsection{Coequalizador}

O coequalizador de $f, g: S_1 \rightrightarrows S_2$ é o quociente de $S_2$ pela menor congruência compatível com $\varphi_2$ que identifica $f(i) \sim g(i)$ para todo $i \in S_1$.

\begin{lstlisting}[style=matlab, caption={Coequalizador em Struct}]
function [CE, quot] = struct_coequalizer(S1, S2, f, g)
    parent = 1:S2.n;
    function r = find_root(x)
        while parent(x) ~= x, x = parent(x); end
        r = x;
    end
    % Identificar f(i) ~ g(i)
    for i = 1:S1.n
        a = find_root(f(i));
        b = find_root(g(i));
        if a ~= b, parent(a) = b; end
    end
    % Fechar para phi2: identificar phi2(f(i)) ~ phi2(g(i))
    changed = true;
    while changed
        changed = false;
        for x = 1:S2.n
            a = find_root(x);
            b = find_root(S2.phi(x));
            pa = find_root(S2.phi(a));
            if pa ~= find_root(b)
                parent(pa) = find_root(b);
                changed = true;
            end
        end
    end
    classes = zeros(1, S2.n);
    rep = 0;
    for x = 1:S2.n
        r = find_root(x);
        if classes(r) == 0
            rep = rep + 1;
            classes(r) = rep;
        end
        classes(x) = classes(find_root(x));
    end
    CE.n   = rep;
    CE.phi = zeros(1, rep);
    CE.d   = zeros(1, rep, 'like', 1i);
    for x = 1:S2.n
        cx = classes(x);
        CE.phi(cx) = classes(S2.phi(x));
        CE.d(cx)   = S2.d(x);
    end
    quot = classes;
end
\end{lstlisting}